\subsection{Ranking Schemes and Stored Output}
\label{sec:ranking}

Rationale Miner ranked its candidates under two schemes, a sentence-based scheme (SBS) that orders sentences by their heuristic score and a message-based scheme (MBS) that orders whole messages, and found that MBS recovered more of the ground truth in the top ranks---61.5\% against 39.7\% of the accepted-PEP rationales in the top five, 74.4\% against 48.0\% for rejected PEPs---because a message that carries several signals is read once, in context \cite{sharma2021rationale}. Influence Miner keeps both schemes, and this section makes explicit what is stored, how each scheme orders it, and which scheme each part of the study uses.

\paragraph{What is stored.} Step~8 writes one row per candidate sentence into \texttt{influence\_candidates} (Table~\ref{tab:fields}). Every label is stored next to the evidence that produced it: the cue phrases that fired, per mechanism, so that a label can be checked against its sentence without re-running the detector; the score, so that any ordering can be reproduced; and the message id, so that the sentence can be opened in its thread. Nothing is stored for a sentence without a mechanism. On the extended Python corpus the table holds 299,192 rows for 764 PEPs; the scores run from the threshold of 1.0 to 5.4 with a mean of 2.13, 60.1\% of the rows score at least 2.0, 11.3\% at least 3.0 and 0.6\% at least 4.0, so the score separates a small head of strong signals---several mechanisms, an explicit stance, a formal role, proximity to the decision---from a long tail of single-mechanism sentences.

\begin{table}[t]
\caption{Fields stored per influence candidate (\texttt{influence\_candidates}).}
\label{tab:fields}
\footnotesize
\begin{tabular}{p{0.30\columnwidth}p{0.62\columnwidth}}
\toprule
Field & Content \\
\midrule
proposal, message\_id, message\_date & the proposal--message pair the sentence came from (a message linked to two proposals yields two rows) \\
author\_name, author\_email, author\_role & the writer after identity resolution, and the role held on that proposal on that date \\
sentence & the cleaned sentence \\
influence\_types, primary\_influence\_type & the mechanisms detected (multi-label), and the first of them \\
evidence\_cues & the cue phrases that fired, per mechanism (``authority: the council; procedural\_control: pep 13'') \\
influence\_direction, influence\_scope, influence\_target & supporting, blocking, revising or neutral; internal, external, mixed or unknown; what is being influenced \\
controversial & 1 if any of the seven controversial mechanisms is present \\
final\_decision, decision\_date, days\_before\_decision, decision\_phase, governance\_era & the proposal's outcome and the sentence's position in its life \\
aligns\_with\_outcome & whether the direction agrees with the outcome (Section~\ref{sec:features}) \\
influence\_score & the additive heuristic score (Section~\ref{sec:scoring}) \\
extraction\_scheme & the detector version that produced the row (\texttt{heuristic-v2}) \\
\bottomrule
\end{tabular}
\end{table}

\paragraph{Sentence-based scheme.} SBS is the native order of the table: within a proposal, candidates are ranked by score, ties broken by date. It is the scheme behind every corpus statistic in Section~\ref{sec:results}, where the unit is the sentence and each row counts once towards its mechanisms, its direction, its writer's role and its phase. The threshold of 1.0 admits every sentence with a mechanism and at least one further signal, so SBS is an ordering of the whole candidate set rather than a selection from it; the annotation sample (Section~\ref{sec:outputs}) draws from its head, its tail and its strata so that the annotation study can tell where along the ranking precision falls off.

\paragraph{Message-based scheme.} MBS folds the candidates of one message into one row of \texttt{influence\_message\_summary}, rebuilt from the candidate table by one grouped query in step~9: the number of candidate sentences, the sum of their scores, the highest score, the union of their mechanisms, the counts of supporting, blocking and revising sentences, the strongest sentence, and the writer, role and outcome. Messages are ranked by the sum by default---a message that argues on several fronts ranks above a one-line remark---with the maximum as the alternative that favours the single strongest sentence. Table~\ref{tab:mbs} gives the scheme's shape on the Python corpus. Half of the proposal-linked message rows carry at least one candidate; of those, a third carry exactly one and 7\% carry ten or more, and it is the latter---position statements, summaries, pronouncements---that MBS is built to surface. For the 607 PEPs with at least fifty candidates, the ten highest-ranked messages (8\% of a PEP's messages on average) hold half of its candidates under the sum and 37\% under the maximum; the two variants agree on about five of their ten messages; and 57\% of a PEP's ten highest-scoring sentences sit inside its ten highest-ranked messages, so the schemes point at the same core of the discussion from different sides. The schemes also differ in whom they put first: the highest-scoring \emph{sentence} of a PEP is most often a core developer's or the BDFL's, while the highest-ranked \emph{message} is most often the BDFL's, and community members and PEP editors---who write long, structured messages---move up under MBS (Table~\ref{tab:mbs}, lower panel). This is the long-message bias the sum carries, and the reason the maximum is kept beside it.

\begin{table}[t]
\caption{Sentence- and message-based ranking on the extended Python corpus.}
\label{tab:mbs}
\footnotesize
\setlength{\tabcolsep}{4pt}
\begin{tabular}{p{0.56\columnwidth}r}
\toprule
Candidate sentences (SBS rows) & 299,192 \\
Proposal--message rows with $\geq 1$ candidate (MBS rows) & 76,087 \\
\quad as a share of the 150,007 proposal-linked message rows & 50.7\% \\
\quad distinct messages & 64,375 \\
Candidates per message: mean & 3.93 \\
\quad exactly 1 / 2 / 3 / 4--9 / $\geq 10$ (\% of messages) & 35.4 / 21.2 / 12.9 / 23.7 / 6.8 \\
Message score: mean sum / mean maximum & 8.37 / 2.41 \\
Messages with a controversial mechanism & 14.1\% \\
\midrule
\multicolumn{2}{l}{\emph{607 PEPs with $\geq 50$ candidates (mean 124 messages, 488 candidates)}} \\
Share of candidates in the top-10 messages: by sum / by maximum & 51.6\% / 36.8\% \\
Top-10 messages shared by the sum and maximum variants & 4.9 of 10 \\
Top-10 SBS sentences inside the top-10 MBS messages & 56.9\% \\
Candidates per top-10 message (mean) / their maximum score & 15.4 / 3.07 \\
\midrule
\multicolumn{2}{l}{\emph{Role of the top-ranked item per PEP (764 PEPs): SBS sentence / MBS message}} \\
Core developer & 232 / 178 \\
BDFL & 187 / 211 \\
Steering Council & 127 / 82 \\
Proposal author & 112 / 75 \\
Community member & 41 / 122 \\
PEP editor & 36 / 84 \\
BDFL delegate & 29 / 12 \\
\bottomrule
\end{tabular}
\end{table}

\paragraph{Which scheme where.} The corpus statistics use SBS because the research questions are about signals, not messages: how often each mechanism is used, by whom, when, and with what outcome. MBS is the reading order: it is what the viewer shows an analyst who wants to know which messages shaped a proposal, and it is the order in which the annotation study will present messages to annotators, following the finding of the rationale study that annotators judge a signal more reliably when they see the whole message. Both orders are reproducible from the stored fields alone, without re-running the detectors.
